\section{\label{sec:rgeqns}Renormalization group equations}

We consider the matrix elements of scalar operators coupled to 
$N$ external, unsmeared scalar fields. We can derive renormalization group 
equations 
for the Wilson coefficients of the OPE by considering the scale dependence of 
suitably chosen Green functions \cite{Collins:1984rdg}. The Green function for 
$N+2$ external scalar 
fields obeys the renormalization group equation
\begin{equation}\label{eq:rg1}
\left[\mu\frac{\mathrm{d}}{\mathrm{d} \mu}+\left(N+2\right)\gamma 
\right]\big\langle \Omega |  
\widetilde{\phi}(p_1) \ldots\widetilde{\phi}(p_{N+2}) | \Omega \big\rangle = 0,
\end{equation}
while the Green function of the renormalized operator $\phi_{\mathrm{R}}^2(0)$ 
coupled 
to $N$ external scalar fields satisfies
\begin{equation}\label{eq:rg2}
\left[\mu\frac{\mathrm{d}}{\mathrm{d} \mu}-\gamma_m + N\gamma
\right]\big\langle \Omega |  \phi_{\mathrm{R}}^2(0)
\widetilde{\phi}(p_1) \ldots\widetilde{\phi}(p_N) | \Omega \big\rangle = 0.
\end{equation}
Here the renormalization group operator for scalar field theory is
\begin{equation}
\mu\frac{\mathrm{d}}{\mathrm{d} \mu} = \mu\frac{\partial}{\partial 
\mu}\bigg|_{\lambda,m_{\mathrm{R}}} + \beta 
\frac{\partial}{\partial \lambda}\bigg|_{\mu,m_{\mathrm{R}}} - \gamma_m 
m_{\mathrm{R}}\frac{\partial}{\partial m_{\mathrm{R}}}\bigg|_{\mu,\lambda}
\end{equation}
and the coefficients are \cite{Kleinert:2001ax}
\begin{align}
\beta^{\overline{MS}}(\lambda) = {} & \mu 
\frac{\mathrm{d}\lambda}{\mathrm{d}\mu} 
=\frac{3\lambda^2}{(16\pi^2)^2}+{\cal O}(\lambda^3), \\
\gamma_m^{\overline{MS}}(\lambda) = {} & 
-\frac{\mu}{2}\frac{\mathrm{d}\log(m_{\mathrm{R}})}{\mathrm{d}\mu}
\nonumber \\
= {} &  -\frac{\lambda}{2(16\pi^2)}+
\frac{5\lambda^2}{12(16\pi^2)^2}+{\cal O}(\lambda^3), \label{eq:gammam}\\
\gamma^{\overline{MS}}(\lambda) = {} & \frac{\mu}{2} 
\frac{\mathrm{d}\log(Z_\phi)}{\mathrm{d}\mu} 
=  \frac{\lambda^2}{12(16\pi^2)^2}+{\cal O}(\lambda^3).
\end{align}
In general these coefficients depend on the mass, the renormalization scale, 
and 
the renormalized coupling constant, but in a mass independent renormalization 
scheme, such as the $\overline{MS}$ scheme, they depend on the 
renormalization scale only through the renormalized coupling constant. These 
functions are known to five loops for the $O(N)$-symmetric theory, given by the 
$N$-multiplet $\Phi = \{\phi_1,\ldots,\phi_N\}$ 
\cite{Derkachov:1997pf,Kleinert:2001ax}.

Returning again to our example of the OPE for the two-point function, Eq.~
\eqref{eq:opeint}, we can derive an renormalization group equation for the 
Wilson coefficient 
$c_{\phi^2}$ by coupling these operators to two external scalar fields and 
using Eqs.~\eqref{eq:rg1} and \eqref{eq:rg2} (for further details, see, for 
example, Ref.~\cite{Collins:1984rdg}):
\begin{equation}
\left[\mu\frac{\mathrm{d}}{\mathrm{d} \mu} 
 + 2\big(\gamma +\gamma_m\big)\right]c_{\phi^2} = 0.
\end{equation}

Just as we might expect, the anomalous dimension of the Wilson coefficient 
$c_{\phi^2}$ is equal to the difference between the anomalous dimension of 
$\phi(x)\phi(0)$ and that of $[\phi^2(0)]_{\mathrm{R}}$.

\subsection{Renormalization group equations for the sOPE Wilson coefficients}

The flow time, $\tau$, introduces a new scale into the problem. In principle 
one can view the flow time as just another external scale, unrelated to the 
renormalization scale $\mu$. In this case, it is natural to modify the 
renormalization group equations to account for the change in the Green 
functions as we change both $\mu$ and $\tau$:
\begin{equation}
\mu\frac{\mathrm{d}}{\mathrm{d} \mu}\rightarrow  
\mu\frac{\mathrm{d}}{\mathrm{d} 
\mu}- 2\tau \frac{\mathrm{d}}{\mathrm{d} \tau}.
\end{equation}
Our choice of differential operator is constrained by the mass dimension of 
each scale, $\mu$ and $\tau$, but is not unique. In particular, one could 
choose 
the operator $\mu\mathrm{d}/\mathrm{d} 
\mu + \tau \mathrm{d}/\mathrm{d} \tau$. This freedom does not affect 
the logic of our discussion nor our conclusion, because alternative conventions 
can be absorbed into the definition of the renormalization parameters and 
anomalous dimensions in, for example, Eq.~(58).

At this stage it is worth commenting on the two scales in the problem, $\mu$ 
and $\tau$. In DIS, the spacetime separation of the corresponding
OPE, $x$, and renormalization scale $\mu$ are, in principle, two distinct 
scales. The spacetime separation is provided by the inverse momentum transfer 
of a particular DIS experiment or set of experiments. The renormalization 
scale, however, is a theoretical choice, and ultimately physical quantities 
should not depend on the renormalization scale. It is generally convenient to 
choose $\mu = 1/x$, but it is not strictly necessary.

The relationship between the flow time and the renormalization scale is 
analogous and these two scales are distinct. For the sOPE, the flow time can be 
considered as simply an external scale, imposed by some particular lattice 
``experiment'', and the renormalization scale is a convenient theoretical 
choice. We will see that it is helpful to tie these scales 
together, to reduce the two-scale problem to a single scale, but this is not 
formally necessary. Therefore, in the following analysis, the flow and 
renormalization scale should be understood as completely independent scales.

Considering again the sOPE for the two-point function, Eq.~
\eqref{eq:sope2pt},
we determine the renormalization group equation for the smeared Wilson 
coefficient, 
$d_{\rho^2}$, by following a procedure analogous to that outlined above. 

We assume that we are working in the small flow-time regime, which allows us 
to relate 
operators at vanishing and nonvanishing flow time 
\cite{Makino:2014taa,Suzuki:2013gza,Luscher:2011bx}:
\begin{equation}\label{eq:smallt}
[\phi^2(0)]_{\mathrm{R}}  =   {\cal 
Z}_{\rho^2}(\tau,\mu)\rho^2(\tau,0) + {\cal O}(\tau).
\end{equation}
This coefficient satisfies
\begin{equation}\label{eq:zgamma}
\mu\frac{\mathrm{d}}{\mathrm{d}\mu}\log({\cal 
Z}_{\rho^2}(\tau,\mu^2)) = 2\gamma_m,
\end{equation}
and to one loop is given by
\begin{equation}
{\cal Z}_{\rho^2}(\tau,\mu^2)= 
1-\frac{\lambda}{32\pi^2}\left[1+\gamma_{\;\mathrm{E}}+\log(2\tau 
\mu^2)\right]+{\cal 
O}(\lambda^2,\tau).
\end{equation}

We apply the renormalization group operator
\begin{equation}
\mu\frac{\mathrm{d}}{\mathrm{d} \mu} 
-2\tau\frac{\mathrm{d}}{\mathrm{d} \tau} + \left(N+2\right)\gamma
\end{equation}
to matrix elements of the operators in Eq.~\eqref{eq:smallt} coupled to 
$N$ 
external scalar fields to obtain
\begin{equation}\label{eq:srge}
\left[\mu\frac{\mathrm{d}}{\mathrm{d} 
\mu} -2\tau\frac{\mathrm{d}}{\mathrm{d} 
\tau} +2(\zeta_{\rho^2}-\gamma)\right]d_{\rho^2} = {\cal O}(\tau).
\end{equation}
Here
\begin{equation}
\zeta_{\rho^2} = \tau\frac{\mathrm{d}}{\mathrm{d}\tau}\log({\cal 
Z}_{\rho^2}(\tau,\mu^2))\label{eq:zetadef}
\end{equation}
is an anomalous dimension associated with the flow-time dependence of the 
operator $\rho^2(\tau,0)$. The renormalization group equation for the 
corresponding
matrix element of $\rho^2(\tau,0)$ coupled to $N$ external fields is 
given by
\begin{align}\label{eq:rhotau}
\bigg[\mu\frac{\mathrm{d}}{\mathrm{d} 
\mu}{} & -2\tau\frac{\mathrm{d}}{\mathrm{d} 
\tau} +2\zeta_{\rho^2}+N\gamma\bigg] \nonumber\\
{} & \times\big\langle \Omega 
|\rho^2(\tau,0)\widetilde{\phi}(p_1)\ldots\widetilde{\phi}(p_N) 
| \Omega \big\rangle = {\cal O}(\tau).
\end{align}
We note that this equation only holds provided the flow time is small compared 
with the momenta of the external particles, which for DIS would be of the order 
of hadronic scales.

If we now demand that the smearing scale, 
$\tau$, and the inverse of the renormalization scale, $\mu$, are 
proportional to each other, \emph{i.e.}, $\tau = b/\mu^2$ with $b$ real, then 
the 
renormalization group 
equation 
becomes
\begin{align}\label{eq:rhotau_mutau}
{} &
\bigg[2\mu\frac{\mathrm{d}}{\mathrm{d} 
\mu} +2\zeta_{\rho^2}+N\gamma\bigg] \nonumber\\
{} &\qquad \times\big\langle \Omega 
|\rho^2\left( b/\mu^2,0\right)\widetilde{\phi}(p_1)\ldots\widetilde{\phi
} (p_N) 
| \Omega \big\rangle = {\cal O}(b).
\end{align}
This renormalization group equation provides the starting point for a 
nonperturbative 
step-scaling method \cite{Monahan:2013lwa,Luscher:1991oxn} that evolves 
nonperturbative matrix elements to a high scale, where they can be combined 
with 
perturbative smeared Wilson coefficients. Here the renormalization group 
equation holds in 
the small flow time limit, or in other words, provided $b\ll 1$. This 
constraint automatically ensures that the flow time is also smaller than any 
hadronic length scales, $\Lambda_{\mathrm{QCD}} \ll \mu^2 \ll 1/\tau$, and is 
generally true for practical step-scaling methods 
\cite{Monahan:2013lwa,Luscher:1991oxn}. We are currently 
investigating this approach in QCD.
